\begin{table}[p]\centering
\caption{Original nonlinear transfer: certificate size relative to available gain}\label{tab:r12-ratios}
\setlength{\tabcolsep}{3.5pt}\begin{tabular}{llrrr}
\toprule
Family & Actor & \shortstack{Mean raw\\gain} & \shortstack{Proposal gap /\\classical gain} & \shortstack{Selected gap /\\classical gain} \\
\midrule
Cycle & Value & 0.008528 & 0.59 & 0.59 \\
Cycle & Value--gradient & 0.009686 & 0.44 & 0.44 \\
Random-edge & Value & 0.012158 & 0.43 & 0.43 \\
Random-edge & Value--gradient & 0.013160 & 0.24 & 0.24 \\
Geometric & Value & 0.002601 & 3.00 & 3.00 \\
Geometric & Value--gradient & 0.003773 & 2.31 & 2.31 \\
Weighted block & Value & -0.097068 & 1340.86 & 23.07 \\
Weighted block & Value--gradient & -0.098095 & 1484.75 & 23.07 \\
\bottomrule
\end{tabular}
\begin{minipage}{.98\textwidth}\small
Ratios are medians across the original two instances and ten training seeds, using the classical certified feasible gain as denominator. The selected gap uses both the proposal and static upper bounds, as in the main-paper gate equation. The weighted-block gate still has a large gap; a no-harm choice is not an accuracy certificate. All forty weighted-block pipelines were subsequently refined by the cached classical solver and independently audited.
\end{minipage}
\end{table}
